% -----------------------------------------------------------
\section{Blessing, Bless}
% -----------------------------------------------------------
	\label{BLESSING}
	
% % ----------------------
% \subsection{See also}
% % ----------------------

% % ----------------------
% \subsection{1828 American Dictionary of the English Language} % *1828
% % ----------------------
	% \label{BLESSING-1828}

	% \begin{compactitem}
		% \item
	% \end{compactitem}

% % ----------------------
% \subsection{Oxford English Dictionary} % *OED
% % ----------------------
	% \label{BLESSING-OED}

	% \begin{compactitem}
		% \item[]
	% \end{compactitem}

% ----------------------
\subsection{Scriptural Usage} % *SCRIPTURES
% ----------------------
	\label{BLESSING-SCRIPTURES}

	% % -----
	% \subsubsection{Old Testament} % *OT
	% % -----
		% \label{BLESSING-OT}
	
		% % --
		% \paragraph{KJV}
		% % --
			% \label{BLESSING-OT-KJV}
	
			% \begin{tabular}{ll}
				% Total occurrences in the OT & 
			% \end{tabular}
			
			% \begin{compactdesc}
				% \item[]
			% \end{compactdesc}
			
		% % --
		% \paragraph{Hebrew Bible}
		% % --
			% \label{BLESSING-HEBREW}
		
			% %
			% \subparagraph{\texorpdfstring{\he{sample word}}{}}
			% %
				% \label{PAGA} % whatever is in step bible without periods
			
				% \begin{compactdesc}
					% \item[HALOT , PDF ] \textbf{}
						% \begin{compactitem}
							% \item[1]
						% \end{compactitem}
					% \item[BDB , PDF ] \textbf{}
						% \begin{compactitem}
							% \item[1]
						% \end{compactitem}
					% \item[DCH PDF ] \textbf{}
						% \begin{compactitem}
							% \item[1]
						% \end{compactitem}
				% \end{compactdesc}
				
				% \begin{compactdesc}
					% \item[Some OT uses] \textbf{}
						% \begin{compactdesc}
							% \item[]
						% \end{compactdesc}
				% \end{compactdesc}
			
		% % --
		% \paragraph{Septuagint}
		% % --
			% \label{BLESSING-LXX}
		
			% %
			% \subparagraph{\texorpdfstring{\gr{sample word}}{}}
			% %
		
	% -----
	\subsubsection{New Testament} % *NT
	% -----
		\label{BLESSING-NT}
	
		% % --
		% \paragraph{KJV}
		% % --
			% \label{BLESSING-NT-KJV}
			
	
			% \begin{tabular}{ll}
				% Total occurrences in the NT & 
			% \end{tabular}
			
			% \begin{compactdesc}
				% \item[]
			% \end{compactdesc}
			
		% --
		\paragraph{Greek New Testament} % *GREEK
		% --
			\label{BLESSING-GREEK}
		
			%
			\subparagraph{\texorpdfstring{\gr{μακάριος}}{}}
			%
				\label{MAKARIOS}
			
				\begin{compactdesc}
					\item[BDAG 541b] \textbf{}
						\begin{compactitem}
							\item[1] pertaining to being fortunate or happy because of circumstances, \textit{fortunate, happy}
							\item[1a] of humans
							\item[1b] of transcendent beings, viewed as \textit{privileged, blessed}
							\item[2] pertaining to being especially favored, \textit{blessed, fortunate, happy, privileged}
							\item[2a] of humans \textit{privileged recipient of divine favor}
							\item[2b] of things or experiences \textit{blessed}
							\item[2b$\alpha$] of parts of the body of persons who are the objects of special grace, which are themselves termed blessed
							\item[2b$\beta$] of things that stand in a very close relationship to the divinity
							\item[2b$\gamma$] of martyrdoms...of the object of Christian hope
						\end{compactitem}
					% \item[LSJ , PDF ] \textbf{}
						% \begin{compactitem}
							% \item 
						% \end{compactitem}
					\item[NIVAC NT 01 291--292] \textit{Makarios} is a state of existence in relationship to God in which a person is ``blessed'' from God's perspective even when he or she doesn't feel happy or isn't presently experiencing good fortune. This does not mean a conferral of blessing or an exhortation to live a life worthy of blessing; rather, it is an acknowledgment that the ones indicated are blessed. Negative feelings, absence of feelings, or adverse conditions cannot take away the blessedness of those who exist in relationship with God.
					\item[WBC 33a 91] Although the word \gr{μακάριοι}, which appears as the first word in each of the nine beatitudes, occurs in Hellenist literature, where it describes those of good fortune, the true background to the NT use of the word is in the OT (Zimmerli finds forty-six instances in the Hebrew canon). The LXX often uses the word as a translation of \he{'a+s:rey} (deeply ``happy, blessed''). The word is of course especially appropriate in the NT in such contexts as the present one, where it describes the nearly incomprehensible happiness of those who participate in the kingdom announced by Jesus. Rather than happiness in its mundane sense, it refers to the deep inner joy of those who have long awaited the salvation promised by God and who now begin to experience its fulfillment. The \gr{μακάριοι} are the deeply or supremely happy.
				\end{compactdesc}
				
				\begin{compactdesc}
					\item[Some NT uses] \textbf{}
						\begin{compactdesc}
							\item[]
						\end{compactdesc}
				\end{compactdesc}
		
	% % -----
	% \subsubsection{Book of Mormon} % *BOM
	% % -----
		% \label{BLESSING-BOM}
	
		% \begin{tabular}{ll}
			% Total occurrences in the BOM & 
		% \end{tabular}
		
		% \begin{compactdesc}
			% \item[]
		% \end{compactdesc}
		
	% % -----
	% \subsubsection{Doctrine and Covenants} % *DC
	% % -----
		% \label{BLESSING-DC}
	
		% \begin{tabular}{ll}
			% Total occurrences in the DC & 
		% \end{tabular}
		
		% \begin{compactdesc}
			% \item[]
		% \end{compactdesc}
		
	% % -----
	% \subsubsection{Pearl of Great Price} % *PGP
	% % -----
		% \label{BLESSING-PGP}
	
		% \begin{tabular}{ll}
			% Total occurrences in the PGP & 
		% \end{tabular}
		
		% \begin{compactdesc}
			% \item[]
		% \end{compactdesc}
		
% % ----------------------
% \subsection{Key Phrases} % *KEYPHRASES *PHRASES
% % ----------------------
	% \label{BLESSING-PHRASES}

	% \begin{compactdesc}
		% \item[] \textbf{\#times} | list
			% % \begin{compactdesc}
				% % \item[Bible anchor verses] \phantom{.}
					% % \begin{compactdesc}
						% % \item[Romans 5:5] \gr{}
					% % \end{compactdesc}
				% % \item[Commentary] ---
			% % \end{compactdesc}
	% \end{compactdesc}
	
% % ----------------------
% \subsection{Bible Dictionaries} % *DICTIONARIES
% % ----------------------
	% \label{BLESSING-BD}

% % ----------------------
% \subsection{Restoration Usage} % *RESTORATION
% % ----------------------
	% \label{BLESSING-RESTORATION}

	% % -----
	% \subsubsection{Joseph Smith} % *JS
	% % -----
		% \label{BLESSING-JS}
	
		% \begin{compactdesc}
			% \item[{\hyperref[KEY]{Date/Source}}]
		% \end{compactdesc}
	
	% % -----
	% \subsubsection{Brigham Young} % *BY
	% % -----
		% \label{BLESSING-BY}
	
		% \begin{compactdesc}
			% \item[{\hyperref[KEY]{Date/Source}}]
		% \end{compactdesc}
	
% % ----------------------
% \subsection{Commentary and Studies} % *COMMENTARY *STUDIES
% % ----------------------
	% \label{BLESSING-COMMENTARY}

	% % -----
	% \subsubsection{Key Points}
	% % -----
		% \label{BLESSING-KEYS}
	
		% \begin{compactitem}
			% \item
		% \end{compactitem}
		
	% % -----
	% \subsubsection{Sample Study}
	% % -----
		% \label{BLESSING-study}